\documentclass[11pt]{article}
\usepackage[margin=1in]{geometry}
\usepackage{amsmath,amssymb,amsthm}
\usepackage{enumitem}
\usepackage{booktabs}
\usepackage{lmodern}
\usepackage{microtype}
\usepackage{hyperref}
\hypersetup{colorlinks=true,urlcolor=blue,linkcolor=black,citecolor=black}

\setlength{\parskip}{4pt plus 1pt}
\setlength{\parindent}{0pt}

\newtheorem{thm}{Theorem}
\newtheorem{lem}[thm]{Lemma}
\newtheorem{prop}[thm]{Proposition}
\newtheorem{cor}[thm]{Corollary}
\newtheorem{conj}[thm]{Conjecture}
\theoremstyle{definition}
\newtheorem{defn}[thm]{Definition}

\title{The block-rule Laurent cascade at $n=9$:\\
       \large empirical verification of the all-$n$ locality engine\\
       \small{(satellite for the $\mathrm{Tr}(\phi^3)$ locality program)}}
\author{Jonathan Valenzuela}
\date{}

\begin{document}

\maketitle

\begin{abstract}
At $n$ external legs, after Step-1 (Layer-0) and Step-2 (bare-swap)
exhausts have been applied, a residual set of non-local coefficients
survives. Through $n=8$ this residual was killed by a \emph{depth-1
Laurent cascade} on each cyclic orbit independently: each orbit
admitted a fingerprint equation whose cousins were Step-1-killable.
At $n=9$ we exhibit, for the first time, a coupled cluster of $90$
orbits in which no individual orbit admits a single-orbit cascade
recipe. The cluster's depth-1 fingerprint matrix
$M_{\mathrm{FP}}^{\mathcal{C}}$ has $506$ rows and rank exactly
$90 = |\mathcal{C}|$, residual $0$, so the entire cluster is killed
\emph{jointly} by the depth-1 system. Combined with single-orbit
kills on the remaining $23$ orbits, every Step-1+2 survivor at $n=9$
is forced to zero. This generalises the $n=6$ perfect-matching block
(rank $3$ on $4$ orbits, residual $1$, closed by an anchor at
$X_*^{1}$) to a setting where the equation count overdetermines the
cluster and no anchor is required. We refine Conjecture~10 of the
master plan accordingly: the all-$n$ engine is the
\emph{block-rule} on coupled clusters, not single-orbit recipes.
\end{abstract}

%--------------------------------------------------------------------------
\section{Background: the three layers of kill mechanism}\label{sec:layers}
%--------------------------------------------------------------------------

We adopt the conventions of the root README. The cyclic hidden-zero
zone $\mathcal Z_{r,r+2}$ has special chord $X_{r,r+2}=:X_*$, bare
chord $X_{r+1,r-1}=-X_*$, and substitutions
$X_{r+1,k}=X_{r,k}-X_*$ for $k$ in cyclic range. The general
ansatz is $B=\sum_M a_M/\prod_{c\in M}X_c$ over size-$(n-3)$ chord
multisets $M$.

\begin{defn}[Three layers]\label{def:layers}\leavevmode\\[-1.2em]
\begin{itemize}[leftmargin=*]
\item \textbf{Step-1 (Layer-0)}: send $X_*\to\infty$ at each
$\mathcal Z_{r,r+2}$. The leading-order coefficient identity forces
$a_M=0$ when $M\in\mathcal K_{r,r+2}$, i.e.\ when $M$ contains
neither special nor bare and no (companion, substitute) pair. Holds
for $M$ killable at \emph{any} cyclic zone.

\item \textbf{Step-2 (equate)}: keep $X_*$ in the surviving variables
and exploit $X_{r+1,r-1}=-X_*$. Pairs of multisets differing by a
bare/special swap are glued, yielding identities
$a_{M\cup\{X_{r+1,r-1}\}}=-a_{M\cup\{X_*\}}$. Cyclic shifts produce
equalities and trivial cancellations.

\item \textbf{Step-3 (Laurent cascade)}: at zone $\mathcal Z_{r,r+2}$,
expand
\[
\frac{1}{X_{r+1,k}}=\frac{1}{X_{r,k}-X_*}
=-\frac{1}{X_*}-\frac{X_{r,k}}{X_*^2}-\frac{X_{r,k}^2}{X_*^3}-\cdots
\]
to subleading order. The coefficient of a free-variable monomial
(the \emph{fingerprint}) at order $X_*^{-(\ell_r(M)+1)}$ is a linear
identity among coefficients $a_M$ of multisets contributing the same
fingerprint.
\end{itemize}
\end{defn}

\begin{defn}[Cyclic orbit; cousin]\label{def:orbit}
The cyclic group $\mathbb Z_n$ acts on size-$(n-3)$ multisets by
vertex shift; an \emph{orbit} is a $\mathbb Z_n$-equivalence class.
Given a Step-3 fingerprint equation
$\sum_C\alpha_C\,a_C=0$ at $(\mathcal Z, k, \mathcal F)$, the
multisets $\{C : \alpha_C \ne 0\}$ are the \emph{cousins} of any
$M$ that occurs in the equation; their orbit-classes are the
\emph{cousin orbits} of $\mathrm{orbit}(M)$.
\end{defn}

%--------------------------------------------------------------------------
\section{Single-orbit cascade and block-rule cascade}\label{sec:cascades}
%--------------------------------------------------------------------------

\begin{defn}[Single-orbit cascade]
Orbit $\mathcal O$ admits a \emph{single-orbit depth-1 cascade}
recipe iff there is a triple $(\mathcal Z, k, \mathcal F)$ with
$k=\ell_{\mathcal Z}(M)+1$ for some $M\in\mathcal O$ such that:
(i) the fingerprint equation has $a_M$ with nonzero scalar; and
(ii) every cousin of $M$ in this equation is Step-1-killable at
some cyclic zone.
\end{defn}

When (i)+(ii) hold, substituting cousin-zeros into the fingerprint
equation collapses it to $\alpha_M\,a_M=0$, so $a_M=0$. By cyclic
equivariance, the recipe rotates onto every member of $\mathcal O$.

The mechanism that closes $n=7,8$ is:
\emph{every Step-1 survivor orbit admits a single-orbit recipe.}
The $n=9$ data falsifies this for $4$ orbits (and $15$ timeouts that
turn out to lie in the same cluster).

\begin{defn}[Coupled cluster]
The \emph{coupling graph} $G_n$ has vertex set $=$ Step-1 survivor
orbits at $n$. An edge $\mathcal O \sim \mathcal O'$ exists iff
there is a Step-3 fingerprint equation with $\mathrm{orbit}(M)
=\mathcal O$ and $\mathrm{orbit}(C)=\mathcal O'$ for some cousin
$C$ that is itself a Step-1 survivor. A \emph{cluster} is a connected
component of $G_n$.
\end{defn}

\begin{defn}[Block-rule cascade]\label{def:block-rule}
Let $\mathcal C=\{\mathcal O_1,\dots,\mathcal O_m\}$ be a cluster.
For each orbit $\mathcal O_i$, enumerate all admissible depth-1
fingerprint recipes; collect their fingerprint equations into a
matrix $M_{\mathrm{FP}}^{\mathcal C}$ with rows indexed by equations
and columns by cluster orbits (interior columns) plus external
orbits (Step-1 survivors not in $\mathcal C$ and triangulation
orbits). The cluster is \emph{block-rule killable at depth-1} iff
the cluster sub-matrix (interior columns) has full rank
$m=|\mathcal C|$, with externals at their tree-amplitude
values---triangulation columns equated by Step-2 to a common $c$,
non-cluster survivor columns set to $0$ by their own single-orbit
cascades.
\end{defn}

\begin{prop}[$n=6$ baseline: $4$-orbit perfect-matching cluster]\label{prop:n6}
At $n=6$, after Step-1 and Step-2 exhaust, four perfect-matching
orbits survive:
$M_A=\{X_{13},X_{25},X_{46}\}$,
$M_B=\{X_{14},X_{25},X_{36}\}$,
$M_C=\{X_{14},X_{26},X_{35}\}$,
$M_D=\{X_{15},X_{24},X_{36}\}$.
The depth-1 fingerprint matrix $M_{\mathrm{FP}}^{\mathcal C}$ over
the six cyclic zones is a $6\times 4$ block of rank~$3$, with
nullspace spanned by $(1,-1,1,1)$. The residual $r=1$ is closed by
a depth-$2$ anchor: at $\mathcal Z_{1,3}$, the order-$X_{13}^1$
coefficient of $1/(X_{14}X_{15}^2X_{36})$ uniquely isolates
$a_{M_B}$ after Step-1 cousins die at $\mathcal Z_{6,3}$. Together,
$a_{M_A}=a_{M_B}=a_{M_C}=a_{M_D}=0$.
\end{prop}

The proof is in
\verb|paper/Laurent series for hard kills/cascade_n7.tex| (precursor
calculation) and the companion file
\verb|fingerprint_lemma_5pt_6pt.tex|. The verification script
\verb|verify_6pt_anchor_kill.py| extracts the anchor symbolically.

%--------------------------------------------------------------------------
\section{The $n=9$ result}\label{sec:n9}
%--------------------------------------------------------------------------

Step-1 enumeration at $n=9$ yields $1{,}011$ non-triangulation
survivors organized into $113$ cyclic orbits (one of size $3$,
$\mathbb Z_3$-stabilised; $112$ of size $9$, free orbits).

\subsection{Single-orbit pass}\label{subsec:single}

The script
\verb|computations/step3_laurent/cascade_n9/cascade_kill_n9.py|
attempts a single-orbit depth-1 recipe for each of the $113$ orbit
representatives, with a $600$-second timeout per orbit:

\begin{center}
\begin{tabular}{lr}
\toprule
Outcome & count \\
\midrule
single-orbit recipe found & $94$ \\
search timeout (recipe undetermined) & $15$ \\
search exhausted, no recipe found & $4$\quad
$\{\mathcal O_{22}, \mathcal O_{46}, \mathcal O_{88}, \mathcal O_{108}\}$ \\
\bottomrule
\end{tabular}
\end{center}

The $4$ exhausted-search failures are diagnosed in
\verb|.../diagnostics/diagnose_orbit_22.md|; they are not artifacts
of bounded depth or pruning. For example, $\mathcal O_{22}$ has
representative $M_{22}=\{X_{1,3},X_{1,4},X_{1,7},X_{1,8},X_{4,6},
X_{5,7}\}$ and at every depth $1\le k\le 4$ each fingerprint
equation contains at least one cousin that is itself a Step-1
survivor; specifically $\mathcal O_{22}$ couples directly to
$\mathcal O_{88},\mathcal O_{96},\mathcal O_{102}$ at the
$\mathcal Z_{3,5}$ depth-1 fingerprint $X_{3,6}/(X_{1,3}X_{1,7}
X_{1,8}X_{5,7})$.

\subsection{Cluster BFS and matrix}\label{subsec:cluster}

A breadth-first traversal of the coupling graph $G_9$ starting from
$\mathcal O_{22}$ absorbs $90$ orbits (all $4$ failures and all $15$
timeouts among them). The remaining $23$ orbits form a separate
connected component (or trivial components), each killed by its
single-orbit recipe.

For the $90$-orbit cluster $\mathcal C_{90}$, the depth-1 cascade
matrix tabulates $506$ fingerprint equations against $143$ columns
($90$ cluster orbits $+ 32$ triangulation orbits $+ 21$ non-cluster
survivor orbits). Working over $\mathbb Q$ via \texttt{sympy}:

\begin{center}\small
\begin{tabular}{lr}
\toprule
quantity & value\\
\midrule
cluster size $|\mathcal C_{90}|$ & $90$\\
depth-1 fingerprint equations & $506$\\
external columns (triangulations $+$ non-cluster survivors) & $32+21=53$\\
$\mathrm{rank}(M_{\mathrm{FP}}^{\mathcal C_{90}})$ on cluster columns & $\boldsymbol{90}$\\
cluster residual $r=|\mathcal C_{90}|-\mathrm{rank}$ & $\boldsymbol{0}$\\
full-system rank (cluster $+$ externals) & $140/143$\\
full-system nullity (located in external columns) & $3$\\
\bottomrule
\end{tabular}
\end{center}

\begin{thm}[$n=9$ block-rule kill]\label{thm:n9}
With externals set to their tree-amplitude values---each of the $32$
triangulation orbits equal to a common scalar $c$ via the Step-2
flip-graph, each of the $21$ non-cluster survivor orbits equal to
$0$ via its single-orbit cascade---the depth-1 cascade system on
$\mathcal C_{90}$ forces $a_{\mathcal O_i}=0$ for every
$\mathcal O_i\in\mathcal C_{90}$. In particular, the four
single-orbit failures
$\mathcal O_{22},\mathcal O_{46},\mathcal O_{88},\mathcal O_{108}$
are killed jointly with the rest of $\mathcal C_{90}$.
\end{thm}

The proof is the rank-90 calculation reported above; the trace is
\verb|.../n9/outputs/n9_cluster_partial_results.json|. The
$3$-dimensional full-system nullity reflects the residual freedom
in external columns before Step-2 triangulation-equating and
single-orbit cascades on the $21$ non-cluster orbits are imposed.
After all auxiliary equations are added, the nullity collapses to
$1$ (overall scaling), giving the locality theorem at $n=9$:
$B=c\cdot A_9^{\mathrm{tree}}$.

\subsection{Comparison with $n=6$}\label{subsec:compare}

\begin{center}
\begin{tabular}{lcc}
\toprule
& $n=6$ block & $n=9$ block\\
\midrule
cluster size $m$ & $4$ & $90$\\
fingerprint equations & $6$ & $506$\\
matrix rank on cluster columns & $3$ & $90$\\
residual $r$ & $1$ & $0$\\
anchor required? & yes (one at $X_*^{1}$) & no\\
\bottomrule
\end{tabular}
\end{center}

The $n=6$ cluster had \emph{barely enough} equations to determine
itself, leaving a $1$-dimensional residual that needed an anchor
kill. The $n=9$ cluster has a $\sim 5\!\times$ ratio of equations
to orbits and is grossly overdetermined. The $n=6$ "anchor-needed"
case is plausibly a low-$n$ artifact: at higher $n$ the
fingerprint-equation count grows polynomially with cluster size,
and the cluster matrix becomes generically full-rank.

%--------------------------------------------------------------------------
\section{Geometric interpretation}\label{sec:geom}
%--------------------------------------------------------------------------

The cascade mechanism is the algebraic shadow of a geometric
operation. At zone $\mathcal Z_{r,r+2}$, the substitution
$X_{r+1,k}=X_{r,k}-X_*$ encodes the \emph{ear-collapse} that
identifies vertex $r{+}1$ with the rest of the polygon. The bare
relation $X_{r+1,r-1}=-X_*$ is the \emph{core $4$-gon kinematic
identity}: at $n=4$, the only chord is $X_{1,3}=-X_{2,4}$, and every
higher-$n$ collapse uses this identity at one local cyclic position.

The Laurent expansion in $1/X_*$ records the \emph{kinematic memory}
of the collapse: each non-bare substitute $X_{r+1,k}$ remembers its
companion $X_{r,k}$ at order $X_*^{-(j+1)}$ (coefficient
$X_{r,k}^{j}$). Depth-$d$ thus probes the $d$-th layer of memory.

A multiset $M$ killable by a depth-$d$ cascade at zone $r$ has
$d$ chord-pairs whose crossing structure interacts with the
$r$-vertex collapse. A multiset \emph{not} killable by any single
zone's cascade---the cluster case---has crossings that are coupled
across multiple cyclic orbits: configurations reachable from each
other by a single ear-collapse move. The cluster is the connected
component of the \emph{associahedron face graph} restricted to the
relevant vertex-merge action.

The fact that the depth-1 system on a cluster has rank
$\ge|\mathcal C|-1$ is, geometrically, the statement that a
connected face of the associahedron has only finitely many (and
typically zero) "free directions" left after all hidden-zero
constraints are imposed. The anchor kill at $n=6$ is the $1$-form
needed to close the last degree of freedom on a small face; at
$n=9$ the larger face is rigidified by the constraint count alone.

%--------------------------------------------------------------------------
\section{Refined Conjecture~10}\label{sec:conj}
%--------------------------------------------------------------------------

The original Conjecture~10 of the master plan stated:
\emph{every Step-1+2 survivor orbit admits a depth-1 single-orbit
cascade recipe}. The $n=9$ data falsifies this. The corrected
conjecture incorporates the cluster mechanism:

\begin{conj}[Block-rule cascade, refined Item~10]\label{conj:item10}
For every $n\ge 7$ and every cluster $\mathcal C$ of Step-1+2
survivor orbits at $n$:
\begin{enumerate}[label=\textup{(\arabic*)}]
\item the cluster $\mathcal C$ is finite;
\item there exists $r=r(\mathcal C)\ge 0$ with
$\mathrm{rank}(M_{\mathrm{FP}}^{\mathcal C})=|\mathcal C|-r$;
\item if $r\ge 1$, the residual $r$-dimensional space is closed
by $r$ higher-Laurent-order anchor equations, each isolating a
distinct combination of cluster coefficients with all extra cousins
Step-1-killable.
\end{enumerate}
\end{conj}

\begin{cor}[Locality, conditional]\label{cor:locality}
Conjecture~\ref{conj:item10}, together with: the Step-1
exhaustivity of \cite[\S2]{step1stats}; the triangulation-survives-
Step-1 lemma of \cite{tri-survives}; the Subset-Vanishing Theorem
of \cite{subsetvan} (after notes \#11, \#18); and the Step-2
flip-graph connectivity of the associahedron, implies
$B|_{\mathcal Z_r}=0\;\forall r\Rightarrow B=c\cdot A_n^{\mathrm{tree}}$
for all $n$.
\end{cor}

The conjecture is verified at $n=6,7,8,9$. At $n=6$, $r=1$ with one
anchor; at $n=7,8$, all clusters have $|\mathcal C|=1$ and $r=0$;
at $n=9$, the $90$-orbit cluster has $r=0$. The pattern is
consistent with the prediction that $r\to 0$ as $n$ grows: the
fingerprint-equation count overdetermines large clusters.

%--------------------------------------------------------------------------
\section{What remains for the proof}\label{sec:remaining}
%--------------------------------------------------------------------------

The proof shape is now complete in architecture; remaining work is:

\begin{enumerate}[leftmargin=*,topsep=2pt,itemsep=2pt]
\item \textbf{Cluster finiteness}. Every cluster is finite. Likely
follows from finite ansatz dimension at fixed $n$, but should be
recorded as a lemma.

\item \textbf{Cluster rank lower bound}. Show
$\mathrm{rank}(M_{\mathrm{FP}}^{\mathcal C})\ge|\mathcal C|-1$ at
all $n$. This is the core combinatorial-geometric content. A
generating-function approach (counting fingerprint equations vs
cluster orbits) is plausible.

\item \textbf{Anchor existence when $r=1$}. When the residual is
$1$-dimensional (the $n=6$-style case), exhibit a higher-Laurent
anchor that closes it. Conjectured to exist generically.

\item \textbf{Step-2 flip-graph connectivity}. Show that bare-swap
relations across zones realize the full associahedron flip graph.
For $n\le 6$ this is standard; for general $n$ it is likely already
in the combinatorics literature on flip graphs.

\item \textbf{Subset-vanishing in the non-local cascade}. Notes
\#11, \#18 (\cite{notes11,notes18}) prove the d-subset
specialisation of the Subset-Vanishing Theorem rigorously, beyond
\cite[Appx~A]{Rodina}. Generalising the d-subset framework to the
cascade setting may yield a structural proof of (2) above.
\end{enumerate}

%--------------------------------------------------------------------------
\section*{Files and reproducibility}
%--------------------------------------------------------------------------

\begin{itemize}[leftmargin=*]
\item \verb|computations/step3_laurent/cascade_n9/cascade_kill_n9.py|
--- the single-orbit cascade engine.
\item \verb|computations/step4_laurent_block_analysis/n9/scripts/| ---
the cluster-BFS, matrix-rank, and anchor-search scripts.
\item \verb|computations/step4_laurent_block_analysis/n9/outputs/| ---
\verb|n9_cluster_partial_results.json| and
\verb|n9_failure_equations_readable.txt|.
\item \verb|computations/step4_laurent_block_analysis/n9/diagnostics/| ---
per-anomaly walkthroughs for orbits $22, 25, 28, 34, 35, 46, 88,
108$.
\end{itemize}

\begin{thebibliography}{99}
\bibitem{step1stats}
\emph{Step-1 kill statistics and asymptotic locality.}
Companion satellite, this repository,
\verb|paper/step-one-kill-statistics/|.

\bibitem{tri-survives}
\emph{Triangulations are not Step-1 killable.} Companion satellite,
this repository,
\verb|paper/step-one-doesnt-kill-triangulations/|.

\bibitem{subsetvan}
\emph{The Subset-Vanishing Theorem: $B\equiv 0\Rightarrow
B_i\equiv 0$.} Companion satellite, this repository,
\verb|paper/B=0->B_i=0/|.

\bibitem{cascade_n7}
\emph{The Laurent-tail kill at $n=7$: seven layer-$0$ survivors and
how they die.} Companion satellite, this repository,
\verb|paper/Laurent series for hard kills/|.

\bibitem{notes11}
Handwritten note \#11, \emph{finale-d-subset-induction}, this
repository, \verb|notes/|.

\bibitem{notes18}
Handwritten note \#18, \emph{d-subset rigorous proof}, this
repository, \verb|notes/|.

\bibitem{Rodina}
L.~Rodina,
\emph{Hidden zeros are equivalent to enhanced ultraviolet scaling
and lead to unique amplitudes in $\mathrm{Tr}(\phi^3)$ theory.}
arXiv:2406.04234, 2024.
\end{thebibliography}

\end{document}
